\documentclass[12pt]{article}
\usepackage{amsmath, amssymb, amsthm, geometry, hyperref}
\geometry{margin=1in}
\title{Problem 318: 2011 Nines}
\author{Project Euler Solution}
\date{}

\newtheorem{theorem}{Theorem}
\newtheorem{lemma}[theorem]{Lemma}
\newtheorem{corollary}[theorem]{Corollary}
\newtheorem{definition}[theorem]{Definition}

\begin{document}
\maketitle

\section{Problem Statement}

For integers $1 \le p < q$ with $p+q \le 2011$, let $C(p,q,n)$ be the count of consecutive nines at the start of the fractional part of $(\sqrt{p}+\sqrt{q})^{2n}$. Let $N(p,q)$ be the minimal $n$ with $C(p,q,n) \ge 2011$. Find $\sum N(p,q)$ over all valid pairs where $N(p,q)$ is finite.

\section{Mathematical Foundation}

\begin{lemma}[Conjugate Pair]
Define $\alpha = (\sqrt{p}+\sqrt{q})^2 = p+q+2\sqrt{pq}$ and $\beta = (\sqrt{p}-\sqrt{q})^2 = p+q-2\sqrt{pq}$. They are roots of:
\[
    x^2 - 2(p+q)x + (p-q)^2 = 0.
\]
\end{lemma}

\begin{proof}
$\alpha + \beta = 2(p+q)$ and $\alpha\beta = (p+q)^2 - 4pq = (p-q)^2$, both integers.
\end{proof}

\begin{theorem}[Integer Power Sums]
For all $n \ge 0$, $S_n := \alpha^n + \beta^n \in \mathbb{Z}^+$.
\end{theorem}

\begin{proof}
$S_n$ satisfies:
\[
    S_0 = 2,\quad S_1 = 2(p+q),\quad S_n = 2(p+q)S_{n-1} - (p-q)^2 S_{n-2}.
\]
By induction, each $S_n \in \mathbb{Z}$. Since $\alpha > 0$ and $\beta \ge 0$, $S_n > 0$.
\end{proof}

\begin{lemma}[Finiteness Condition]
$N(p,q)$ is finite iff $0 < \beta < 1$ (equivalently $q < (\sqrt{p}+1)^2$) and $pq$ is not a perfect square.
\end{lemma}

\begin{proof}
If $pq$ is a perfect square, $\alpha,\beta \in \mathbb{Z}$ and $\alpha^n$ is an integer. If $pq$ is not a perfect square and $0 < \beta < 1$, then $\{\alpha^n\} = 1 - \beta^n \to 1$, giving arbitrarily many nines. If $\beta \ge 1$ or $\beta = 0$, the fractional part does not approach 1.
\end{proof}

\begin{theorem}[Counting Consecutive Nines]
When $0 < \beta < 1$ and $\beta \notin \mathbb{Z}$:
\[
    C(p,q,n) = \lfloor -n\log_{10}\beta \rfloor.
\]
\end{theorem}

\begin{proof}
$\{\alpha^n\} = 1 - \beta^n$, so the count of leading nines is $\lfloor -\log_{10}\beta^n \rfloor = \lfloor -n\log_{10}\beta \rfloor$.
\end{proof}

\begin{corollary}[Formula for $N(p,q)$]
\[
    N(p,q) = \left\lceil \frac{2011}{-\log_{10}\beta} \right\rceil = \left\lceil \frac{-2011}{\log_{10}(p+q-2\sqrt{pq})} \right\rceil.
\]
\end{corollary}

\begin{proof}
$C(p,q,n) \ge 2011$ iff $n \ge 2011/(-\log_{10}\beta)$. The smallest such integer is the ceiling.
\end{proof}

\section{Algorithm}

\begin{enumerate}
    \item For $p = 1$ to 1005, for $q = p+1$ to $\min(2011-p, \lfloor(\sqrt{p}+1)^2\rfloor)$:
    \begin{itemize}
        \item Skip if $pq$ is a perfect square.
        \item Compute $\beta = p+q-2\sqrt{pq}$; if $0 < \beta < 1$, compute $N = \lceil -2011/\log_{10}\beta \rceil$.
        \item Accumulate $N$.
    \end{itemize}
\end{enumerate}

\section{Complexity Analysis}

\begin{itemize}
    \item \textbf{Time:} $O(P^{3/2})$ where $P \le 2011$, approximately $O(10^{4.5})$.
    \item \textbf{Space:} $O(1)$.
\end{itemize}

\section{Answer}

\[
\boxed{709313889}
\]

\end{document}
